\documentclass[10pt,a4paper]{amsart}
\usepackage{amsmath,amssymb,amsthm,graphicx}
\usepackage{verbatim}
\newenvironment{reference}{\comment}{\endcomment}
%\newenvironment{reference}{}{}
\newenvironment{slogan}{\comment}{\endcomment}
\newenvironment{history}{\comment}{\endcomment}

% For commutative diagrams we use Xy-pic
\usepackage[all]{xy}

% We use 2cell for 2-commutative diagrams.
\xyoption{2cell}
% Theorem environments.
%
\theoremstyle{plain}
\newtheorem{theorem}[subsection]{Theorem}
\newtheorem{proposition}[subsection]{Proposition}
\newtheorem{lemma}[subsection]{Lemma}

\theoremstyle{definition}
\newtheorem{definition}[subsection]{Definition}
\newtheorem{example}[subsection]{Example}
\newtheorem{exercise}[subsection]{Exercise}
\newtheorem{situation}[subsection]{Situation}

\theoremstyle{remark}
\newtheorem{remark}[subsection]{Remark}
\newtheorem{remarks}[subsection]{Remarks}

\numberwithin{equation}{subsection}

% Macros
%
\def\lim{\mathop{\mathrm{lim}}\nolimits}
\def\colim{\mathop{\mathrm{colim}}\nolimits}
\def\Spec{\mathop{\mathrm{Spec}}}
\def\Hom{\mathop{\mathrm{Hom}}\nolimits}
\def\Ext{\mathop{\mathrm{Ext}}\nolimits}
\def\SheafHom{\mathop{\mathcal{H}\!\mathit{om}}\nolimits}
\def\SheafExt{\mathop{\mathcal{E}\!\mathit{xt}}\nolimits}
\def\Sch{\mathit{Sch}}
\def\Mor{\mathop{\mathrm{Mor}}\nolimits}
\def\Ob{\mathop{\mathrm{Ob}}\nolimits}
\def\Sh{\mathop{\mathit{Sh}}\nolimits}
\def\NL{\mathop{N\!L}\nolimits}
\def\CH{\mathop{\mathrm{CH}}\nolimits}
\def\proetale{{pro\text{-}\acute{e}tale}}
\def\etale{{\acute{e}tale}}
\def\QCoh{\mathit{QCoh}}
\def\Ker{\mathop{\mathrm{Ker}}}
\def\Im{\mathop{\mathrm{Im}}}
\def\Coker{\mathop{\mathrm{Coker}}}
\def\Coim{\mathop{\mathrm{Coim}}}

% Boxtimes
%
\DeclareMathSymbol{\boxtimes}{\mathbin}{AMSa}{"02}

%
% Macros for moduli stacks/spaces
%
\def\QCohstack{\mathcal{QC}\!\mathit{oh}}
\def\Cohstack{\mathcal{C}\!\mathit{oh}}
\def\Spacesstack{\mathcal{S}\!\mathit{paces}}
\def\Quotfunctor{\mathrm{Quot}}
\def\Hilbfunctor{\mathrm{Hilb}}
\def\Curvesstack{\mathcal{C}\!\mathit{urves}}
\def\Polarizedstack{\mathcal{P}\!\mathit{olarized}}
\def\Complexesstack{\mathcal{C}\!\mathit{omplexes}}
% \Pic is the operator that assigns to X its picard group, usage \Pic(X)
% \Picardstack_{X/B} denotes the Picard stack of X over B
% \Picardfunctor_{X/B} denotes the Picard functor of X over B
\def\Pic{\mathop{\mathrm{Pic}}\nolimits}
\def\Picardstack{\mathcal{P}\!\mathit{ic}}
\def\Picardfunctor{\mathrm{Pic}}
\def\Deformationcategory{\mathcal{D}\!\mathit{ef}}
\usepackage{cleveref}
\crefname{originalreference}{original reference}{original references}
\Crefname{originalreference}{Original reference}{Original references}
\makeatletter
\newlabel{sites-definition-empty}{{04B5}{source}{}{sourcepointer-1}{}}
\newlabel{sites-definition-empty@cref}{{[originalreference][0][]04B5}{[1][0][]source}{}{}{}}
\makeatother
\crefname{originallemma}{original lemma}{original lemmas}
\Crefname{originallemma}{Original lemma}{Original lemmas}
\makeatletter
\newlabel{lemma-constant-exact}{{093J}{source}{}{sourcepointer-2}{}}
\newlabel{lemma-constant-exact@cref}{{[originallemma][0][]093J}{[1][0][]source}{}{}{}}
\makeatother
\crefname{originallemma}{original lemma}{original lemmas}
\Crefname{originallemma}{Original lemma}{Original lemmas}
\makeatletter
\newlabel{lemma-kernel-surjection-finite-onto-finite-presentation}{{082T}{source}{}{sourcepointer-3}{}}
\newlabel{lemma-kernel-surjection-finite-onto-finite-presentation@cref}{{[originallemma][0][]082T}{[1][0][]source}{}{}{}}
\makeatother
\title{Finiteness and finite presentation of the constant sheaf attached to a module over a site}
\author{The Stacks Project Authors}
\date{stacks-project:sites-modules.tex:local-snapshot-20261012; GNU FDL 1.2 or later}
\begin{document}
\maketitle
\noindent\textit{Editorial note.} All delivered mathematical text is the author's own, copied verbatim from the pinned Stacks Project chapter file sites-modules.tex; no statement or proof step was rewritten, reordered or removed.\par\medskip
\noindent\textit{Editorial note.} The counted span cites one sites-chapter definition, recorded in assembly.external\_references with its canonical tag as a navigation pointer.\par\medskip
\noindent\textit{Editorial note.} The chapter body uses the project's own reference/slogan/history display environments; the retained two-slice preamble of this source defines them verbatim, so the item compiles as in the source.\par\medskip
\noindent\textit{Allowed standard prerequisites (editorial).}\begin{itemize}
\item Sheaf theoretically empty objects and the empty topos. Stacks Project, Sites, Definition sites-definition-empty (tag 00Y0). Supplies the notion of a nonempty object used to reduce the converse direction to a single test object.
\end{itemize}\medskip
\ifdefined\hypertarget\hypertarget{sourcepointer-1}{}\fi
\noindent\textit{Source reference.} Original reference 04B5: sources/stacks-chow/tags.txt, source lines 5640--5640. Cross-chapter standard prerequisite cited by the author inside the counted span; it is used as a pointer and its own proof is not part of this item.\par\medskip
\ifdefined\hypertarget\hypertarget{sourcepointer-2}{}\fi
\noindent\textit{Source reference.} Original lemma 093J: sources/stacks-chow/tags.txt, source lines 11517--11517. Cross-chapter standard prerequisite cited by the author inside the counted span; it is used as a pointer and its own proof is not part of this item.\par\medskip
\ifdefined\hypertarget\hypertarget{sourcepointer-3}{}\fi
\noindent\textit{Source reference.} Original lemma 082T: sources/stacks-chow/tags.txt, source lines 10263--10263. Cross-chapter standard prerequisite cited by the author inside the counted span; it is used as a pointer and its own proof is not part of this item.\par\medskip
\begin{lemma}
\label{lemma-locally-constant-finite-type}
Let $\mathcal{C}$ be a site. Let $\Lambda$ be a ring and let $M$ be a
$\Lambda$-module. Assume $\Sh(\mathcal{C})$ is not the empty topos. Then
\begin{enumerate}
\item $\underline{M}$ is a finite type sheaf of $\underline{\Lambda}$-modules
if and only if $M$ is a finite $\Lambda$-module, and
\item $\underline{M}$ is a finitely presented sheaf of
$\underline{\Lambda}$-modules if and only if $M$ is a 
finitely presented $\Lambda$-module.
\end{enumerate}
\end{lemma}

\begin{proof}
Proof of (1). If $M$ is generated by $x_1, \ldots, x_r$ then
$x_1, \ldots, x_r$ define global sections of $\underline{M}$ which
generate it, hence $\underline{M}$ is of finite type. Conversely, assume
$\underline{M}$ is of finite type. Let $U \in \mathcal{C}$ be an object
which is not sheaf theoretically empty
(Sites, Definition \ref{sites-definition-empty}).
Such an object exists as we assumed $\Sh(\mathcal{C})$ is not
the empty topos. Then there exists a covering $\{U_i \to U\}$
and finitely many sections $s_{ij} \in \underline{M}(U_i)$ generating
$\underline{M}|_{U_i}$. After refining the covering we may assume
that $s_{ij}$ come from elements $x_{ij}$ of $M$. Then
$x_{ij}$ define global sections of $\underline{M}$ whose restriction
to $U$ generate $\underline{M}$.

\medskip\noindent
Assume there exist elements $x_1, \ldots, x_r$ of $M$ which define
global sections of $\underline{M}$ generating $\underline{M}$
as a sheaf of $\underline{\Lambda}$-modules. We will show that
$x_1, \ldots, x_r$ generate $M$ as a $\Lambda$-module.
Let $x \in M$. We can find a covering $\{U_i \to U\}_{i \in I}$
and $f_{i, j} \in \underline{\Lambda}(U_i)$ such that
$x|_{U_i} = \sum f_{i, j} x_j|_{U_i}$. After refining the covering
we may assume $f_{i, j} \in \Lambda$. Since $U$ is not
sheaf theoretically empty, there is at least one $i \in I$
such that $U_i$ is not sheaf theoretically empty.
Then the map $M \to \underline{M}(U_i)$ is injective (details omitted).
We conclude that $x = \sum f_{i, j}x_j$ in $M$ as desired.

\medskip\noindent
Proof of (2). Assume $\underline{M}$ is a $\underline{\Lambda}$-module
of finite presentation. By (1) we see that $M$ is of finite type.
Choose generators $x_1, \ldots, x_r$ of $M$ as a $\Lambda$-module.
This determines a short exact sequence
$0 \to K \to \Lambda^{\oplus r} \to M \to 0$ which
turns into a short exact sequence
$$
0 \to \underline{K} \to \underline{\Lambda}^{\oplus r} \to \underline{M} \to 0
$$
by Lemma \ref{lemma-constant-exact}. By
Lemma \ref{lemma-kernel-surjection-finite-onto-finite-presentation}
we see that $\underline{K}$ is of finite type. Hence $K$ is a
finite $\Lambda$-module by (1). Thus $M$ is a $\Lambda$-module
of finite presentation.
\end{proof}

\end{document}
